\documentclass[]{article}
\usepackage[hidelinks]{hyperref}
\usepackage{graphicx}
\usepackage{amsmath}
%for the bibliography
\usepackage[backend=biber,style=numeric]{biblatex}
\addbibresource{Sources.bib}

\newcommand{\commandToText}[1]{\texttt{\string#1}}

%opening
\title{A \LaTeX  \ Showcase}
\author{Fabrice Kirsch}

\begin{document}

\maketitle

\begin{abstract}
  Today we will learn about \LaTeX.\\
  This is the abstract. Here you put a short summary of your report.
  The tradtional structure of a scientific work is:
  \begin{enumerate}
    \item Introduction
          \begin{itemize}
            \item The current state of the art.
            \item Everything your reader should know.
            \item End with the problem Statement
          \end{itemize}
    \item Materials and Methods (Procedure)
    \item Results
    \item Discussion
          \begin{itemize}
            \item How does your findings fit the state of the art?
          \end{itemize}
    \item (Outlook)
  \end{enumerate}
\end{abstract}

\newpage
\tableofcontents



\section{(Abbreviated) Basics of Scientific Writing}
Main guidelines:
\begin{enumerate}
  \item The content is the interesting part, not the formatting.
  \item Plagiarism: The reader assumes everything you write is your own creation/thought. If it is not cite it.
  \item Add in your mind "unless your supervisor says something different" to the end of every sentence today.
\end{enumerate}
More in the individual sections: \autoref{sec:fig}, \autoref{sec:tab} \\
\noindent For your \textbf{thesis} read: A Manual for Writers of Research Papers, Theses, and
Dissertations by Kate L. Turabian (Library code: 00/ALC 16-7 ) \cite{turabian} \\


\newpage

\section{What is \LaTeX \ and why use it?}
Latex is a text document creation framework. It looks and feels a bit like writing computer code. \\
Advantages:
\begin{itemize}
  \item Looks good out of the box
  \item Algorithmic layout decisions
  \item Formulas look nice
  \item Active helping community (If you cannot google it, I don't know it either)
  \item Some error correction
  \item Makes you feel like a hacker =)
\end{itemize}
Disadvantages:
\begin{itemize}
  \item Takes time to get used to e.g.
        \begin{itemize}
          \item \commandToText{\\} ends a line
          \item Double spaces are ignored
        \end{itemize}
\end{itemize}

\noindent Usually 3 Documents:
\begin{itemize}
  \item .tex file (main document)
  \item .bib file (Bibliography)
  \item .pdf file
\end{itemize}
... and all the figures as separate files.\\
The compiler creates the pdf from the other documents.

(Show Thesis (early version))

\subsection{A few Words Regarding Overleaf}
\begin{itemize}
  \item Web access.
  \item Most packages already installed.
  \item Free individual version.
  \item Free collaborative version is included in Sciebo.
  \item In the pro version you can invite collaborators.
  \item I like TeX Studio a bit better. (Also free, but local)
\end{itemize}

\section{Document Structure}
\begin{figure}[h]
  \centering
  \includegraphics[width=0.7\textwidth]{Structure.jpg}
  \caption{A guidline for students on the structure of their thesis \cite{Singh.2014}}
  \label{fig:structure}
\end{figure}
\textbf{Fabrice include the use case for Citavi here (\autoref{fig:structure})!!!}\\
Avoid having two subsequent headers without text in between.
\section*{Section}
\subsection*{Subsection}
\subsubsection*{Subsubsection}

\section{Figures}\label{sec:fig}
A figure needs a caption below the figure. The caption has to work without the text. The text has to work without the figure. A cross reference is required at a fitting position in the text (\autoref{fig:foto}). Avoid formulations like: "In \autoref{fig:foto} we see..." or "\autoref{fig:foto} shows...". Just write what is seen in the figure.
Figures and tables are floating objects, LaTeX will decide where to place it. Adding a [h!] to the command will try to force to the position in the .tex document.
\begin{figure}[h!]
  \centering
  \includegraphics[width=0.4\textwidth]{Structure.jpg}
  \caption[Short Name for Table of figures]{An example figure showing my face. Just to clarify I would have to credit my wife for the photo because I did not do it myself.}
  \label{fig:foto}
\end{figure}


\section{Tables}\label{sec:tab}
Tables work the same as figures. Only the caption is on top of the table (\autoref{tab:tab}).
A good table only has three lines:
\begin{enumerate}
  \item One above the first row.
  \item One below the first row.
  \item One below the last row.
\end{enumerate}

\begin{table}[h!]
  \centering
  \caption{Caption}
  \begin{tabular}{|c| c| c| c|}
    \hline
    header  1 & header 2 & header 3 & header 4 \\
    \hline
    1         & 2        & 6        & 7        \\
    \hline
    3         & 4        & 9        & 7        \\
    \hline
  \end{tabular}

  \label{tab:tab}
\end{table}



\section{Equations}
Equations c\"an be $\sigma      $ included in different ways. In the text using single\$ signs like $f(x)=\left(x_{2x}\right)$.
Or in a single line using double \$\$ signs
$$
  f(x)=e^{-\frac{x^2}{\sigma}}
$$
$$
  f(x)=e^{-\dfrac{x^2}{\sigma}}
$$
Or we use the equation environment:
\begin{equation}
  a_{index}
\end{equation}
\begin{equation}\label{eq:eq}
  f(x)=e^{-\dfrac{x^2}{\sigma}}
\end{equation}
There are other environments for aligned equal signs
\begin{align}
  f(x) & = x^2+3 \\
       & =5
\end{align}
Of course we can cross reference (\autoref{eq:eq}) or just (\ref{eq:eq})
\section{Citations}
We need a .bib file \cite{Wouters.2022} and in the end we add commands \commandToText{\cite\{name\}} \\
References only appear in the Reference list is they are used. \cite{Glover.1994}
An easy way to get the bib for a source is google scholar
\newpage
\noindent What is still left to be implemented is the table of \textbf{figures} and the \textbf{table} of tables.
\printbibliography
\addcontentsline{toc}{section}{References}

\end{document}


